\documentclass[11pt]{article}
\usepackage[T1]{fontenc}
\usepackage[utf8]{inputenc}
\usepackage{lmodern,microtype}
\usepackage{amsmath,amssymb}
\usepackage[margin=1in]{geometry}
\usepackage{enumitem}
\usepackage[hidelinks]{hyperref}
\hypersetup{
  pdftitle={Research development summary: compact trace rigidity},
  pdfauthor={},
  pdfsubject={Abridged retrospective summary based on project records},
  pdfcreator={LaTeX},
  pdfproducer={pdfTeX}
}
\ifdefined\pdftrailerid
  \pdftrailerid{}
\fi
\setlength{\emergencystretch}{2em}
\setlist{leftmargin=*,itemsep=0.25em}
\title{Research development summary\\
  \large Trace Sparsity and Arithmeticity for Compact Hyperbolic Surfaces}
\author{}
\date{23 October 2026}

\begin{document}
\maketitle

\noindent\textbf{Discussions and model prompting:}
Alexander Kolpakov (UATX) and Nikolay Bogachev (University of Toronto).

\section{Scope and available records}

This is an abridged retrospective account of the development recorded in
the source project's provenance ledger, audit log, and mathematical notes.
It accompanies the focused compact trace-sparsity manuscript. The summary
reports that recorded account; it does not independently reconstruct a
chronology from original messages.

The principal records are
\nolinkurl{PROVENANCE_INFLUENCE_LEDGER.md},
\nolinkurl{AUDIT_LOG.md},
\nolinkurl{CRITICAL_PATH_AUDIT.md}, and
\nolinkurl{FINAL_ARITHMETICITY_DEDUCTION.md}.
Supporting explanations appear in
\nolinkurl{SPECTRAL_COUNTING_BRIDGE.md} and
\nolinkurl{LOCAL_SQUARE_DENSITY_PROOF.md}.
The accompanying provenance document identifies their source revision.

No original conversation export was available in the audited project
tree. This summary therefore contains no verbatim model quotations,
reconstructed dialogue, or authenticated original prompt. The available
release records do not establish a model version, runtime, token count,
or compute budget. The development described here belongs to an external
model-assisted project.

\section{Problem brief}

The manuscript studies the signed trace set of a cocompact Fuchsian lattice
and its counting function
\[
 N_\Gamma(X)=\#\bigl(T(\Gamma)\cap[-X,X]\bigr).
\]
Its main claim is arithmeticity under the critical bound
\[
 \limsup_{X\to\infty}\frac{\log N_\Gamma(X)}{\log X}\leq 1.
\]
Literal linear growth implies this bound, while bounded clustering implies
linear growth by covering a long interval with bounded-length intervals.
The compact Schmutz and Sarnak conclusions are consequences of the stated
critical-growth theorem. This is a description of the submitted
manuscript's claim, not a report of independent mathematical certification.

\section{Development described by the project records}

\subsection*{Initial framework and approaches that did not close the argument}

The provenance ledger identifies signed traces, the subgroup generated by
squares, trace fields, quaternion algebras, and the trace-of-powers
recurrence as inherited structural ingredients. It describes projective
square-root calculations as an early source of two-trace identities.

The records then distinguish several limitations. Difference-set density
can produce arbitrarily close pairs without contradicting bounded
clustering. Scalar gradings and additive-energy arguments do not by
themselves recover the needed arithmetic structure. A route through a
universal radical lower bound for one fixed recurrence would entail
strong unresolved statements of Mersenne and Lucas--Wieferich type.
These are recorded obstructions to particular approaches, not additional
assumptions of the focused theorem.

\subsection*{The moving-element change of strategy}

The ledger identifies fixing a Chebyshev recurrence index \(N\) while
varying the group element as the decisive change. The manuscript uses
elements of the form \(a_m=b^m h_m\), with a short perturbation \(h_m\).
The aim is to retain the relevant local expansion while keeping the
prime-power depths of the recurrence ideals bounded.

The documented argument combines four stages:
\begin{enumerate}
\item A reduced-divisor application of Roth's theorem supplies a lower
  bound with coefficient \(2N-4-\epsilon\).
\item A local prime-square density estimate and bounded-power expansion
  are used in the selection of the varying elements.
\item A uniform two-trace fiber bound controls the multiplicity with
  which group elements represent a trace-graph edge.
\item Pruning the trace graph, congruence counting, and an injection for
  length-two paths supply the complementary upper estimate.
\end{enumerate}

The notes emphasize the parameter order: first let \(m\) tend to infinity,
then let the auxiliary errors tend to zero, and finally let \(N\) tend to
infinity. The claimed resulting estimate compares the local expansion
functional at an algebraic specialization \(y\) with length at the
original character \(x\):
\[
 \mathcal J_{k_y}(\operatorname{tr}_y b)\leq\frac{\ell_x(b)}{2}.
\]
The manuscript contains the mathematical statements, uniformity ranges,
and proofs used in these stages.

\subsection*{Arithmeticity and specialization}

At an algebraic character, the distinguished real place already
contributes half the translation length. The manuscript argues that the
upper estimate forces all other nonnegative local contributions to
vanish. Its final deduction applies Takeuchi's criterion to the
finite-index subgroup generated by squares and then uses
commensurability.

For a transcendental character, the records describe algebraic Fuchsian
specializations preserving the relevant trace collisions. The claimed
local estimate gives marked-length domination; Thurston's separating
asymmetric metric is then used to force equality of the characters.
The passage to compact orbifolds uses a torsion-free finite-index
subgroup.

\section{Checks and repairs recorded in the project}

The audit log records repeated mathematical, reference, exposition, and
formal-scope checks. This summary reports those entries without
establishing that they constitute independent human peer review.
Examples of the recorded repairs include:
\begin{itemize}
\item retaining degree factors when the field used in the local expansion
  functional changes;
\item separating the zero-norm branch of the two-trace calculation from
  the quadratic divisor argument;
\item displaying the modulus-versus-radius condition needed to absorb the
  congruence-counting error and choosing constants in the stated order;
\item including finite covolume and finite index in the Takeuchi
  deduction; and
\item separating literal linear growth from the weaker critical-growth
  hypothesis and checking the exact order of the limits.
\end{itemize}

\section{What the accompanying Lean project establishes}

The Lean package proves selected recurrence and polynomial identities,
finite combinatorial inequalities, and limiting coefficient statements.
It also checks conditional deductions from explicitly supplied analytic,
geometric, specialization, and arithmetic hypotheses.

The package does not construct the relevant hyperbolic surfaces or
formalize all of the manuscript's external inputs. Its public
arithmeticity implications use abstract hypotheses. Finite consistency
examples and counterexamples audit those hypotheses and their numerical
translation; they are not geometric realizations. A successful build or
an axiom report containing only standard logical principles does not
by itself establish that all the geometric hypotheses have been supplied.
The manuscript map gives the declaration-level boundary.

\section{Account metadata and publication boundary}

This edited summary omits account handles, account and session
identifiers, conversation links, local filesystem paths, credentials,
export headers, and unrelated dialogue. The contributor names and
affiliations on the first page are intentional public attribution.
Source-repository attribution is recorded separately in the contribution's
provenance; the surrounding public repository identifies its owner.

Raw conversation exports are not included. Auditing any later transcript
excerpts would require reviewing those actual records and their hidden
metadata. The document above is a summary of available project records,
not an assertion that an unavailable conversation archive has been
anonymized or authenticated.

\end{document}
